\scanpage{378}
\mspage{490}

\[
(29)\quad
\left\{\begin{aligned}
\underset{\substack{\Vert\,\mathrm{dfn.}\\\underline{u}\rho\underline{u}^{-1}}}{\underline{u}(\rho)}
 &=\underset{\substack{\Vert\\\operatorname{int}(\alpha)\rho}}{\underbrace{\xi\rho}},
 &\underline{u}(\sigma)&=\underset{\substack{\Vert\\\varepsilon\operatorname{int}(\beta)(\sigma)}}{\underbrace{\eta\sigma}},\\[1ex]
\det\underline{u}&=\mu
\end{aligned}\right.
\qquad\text{with }\xi=\alpha\rho(\alpha)^{-1},\quad
\eta=\varepsilon\beta\sigma(\beta)^{-1}
\]
\margin{watch the factor $\varepsilon$!}

\uline{Proposition} The morphism of schemes $\Theta$ is
\uline{characterized} by formulas (29), modulo a morphism
$\mathcal{G}_{\rho,\sigma}^{*}\xrightarrow{\lambda}\mu_2$.

Indeed, the first two formulas, which describe the action of
$\underline{u}$ on $\rho$ resp. $\sigma$, determine $\underline{u}$ modulo
$T_\rho$ resp. modulo $T_\sigma$, hence \inserted{(together)} modulo
$T_\rho\cap T_\sigma=\mathbb{G}_{\mathrm m}$, the third determines it
modulo $SG$, -- hence with the preceding ones, modulo $\mu_2$ -- i.e. another
morphism $\Theta'$ having the same properties
will be of the form
\[
\Theta'(g)=\lambda(g)\Theta(g)\qquad\text{where }\lambda:
\mathcal{G}_{\rho,\sigma}^{*}\longrightarrow\unsure{\mathcal{E}_2|\varepsilon_2}.
\]
\unsure{q.e.d.}

\uline{Remark} Let us note \inserted{moreover} that \deleted{\unsure{every}}
$\Theta$ is normalized by
\[
\Theta(\underbrace{1,1,1,1}_{1})=1
\]
which imposes \deleted{\unsure{an}}
\deleted{\unsure{additional}}\inserted{corresponding}
normalization $\lambda(1)=1$; on $\lambda$ -- but \unsure{unless I am mistaken} this
already implies (at least if $\rho,\sigma$ are $\frac12$ simple hence
$T_\rho,T_\sigma$ tori\ldots) that $\lambda$ is a \uline{homomorphism}
of group schemes, which here is \unsure{necessarily} trivial,
on \inserted{the component} $B'_0\times T_\rho\times T_\sigma$
(at least in characteristic $\ne2$) so we find essentially as the only
possible indeterminacy the multiplication \deleted{by} of
$\underline{u}$ by $\varepsilon$ \deleted{\unsure{all}}
(corresponding to the \inserted{canonical} homomorphism (nontrivial) from
$N'_\sigma/T_\sigma$ onto $\mu_2$)\ldots\ \deleted{\unsure{Error}}

\margin{If we worked \unsure{with} $S\mathcal{G}_{\rho,\sigma}$
instead of $\mathcal{G}_{\rho,\sigma}^{*}$, the indeterminacy would be
given by multiplication \unsure{by|$\mapsto$}
$\varepsilon^{i}\varepsilon'^{j}$
$(i,j\in\mathbb{Z}/2\mathbb{Z})$.}
\margin{377}

% Printed archive footer (not autograph).
\begin{center}
377\\
\texttt{https://grothendieck.umontpellier.fr}\\
Université de Montpellier, UMR CNRS 5149, Case courrier 051, Place Eugène Bataillon\\
34095 Montpellier cedex 5 - France
\end{center}

\scanpage{379}
\mspage{\unsure{490}}

(IX) \uline{The groups} $\mathcal{M}_{\rho,\sigma}$ \uline{and}
$S\mathcal{M}_{\rho,\sigma}$.

Let $\rho,\sigma$ be as always two sections of $G$, satisfying the
general conditions (cf. for example proposition page \unsure{478'|470'}), plus the
condition $\operatorname{Tr}\sigma=0$ (cf. \uline{V} -- for example lemma
page 455 and its corollary page 457) -- we could carry out their construction
also without this additional condition (contenting ourselves with
working with $\mathcal{G}_{\rho,\sigma}$ instead of
$\mathcal{G}_{\rho,\sigma}^{*}$) -- but \unsure{it matters little}
probably for what follows. We can if need be suppose
$\rho,\sigma$ normalized so that we have
$\operatorname{Tr}(\rho+\sigma)=1$.

We denote by $S\mathcal{M}_{\rho,\sigma}$ the product scheme
\[
(1)\qquad S\mathcal{M}_{\rho,\sigma}
=S\mathcal{G}_{\rho,\sigma}^{*}\times_S SG
\]
\margin{\unsure{In|For} what follows, the ``\unsure{officialization}''
\illegible{1} of the group \unsure{schemes|structures} $S_0\mathcal{M}_{\rho,\sigma}$,
cf. page 559 ff.}

Its sections are denoted
\[
(\alpha,\beta,\mu,\varepsilon,f)\ ;
\]
$f$ thus being an arbitrary section of $SG$
$(\simeq\mathrm{Sl}(2)_S)$.

In the case where $S=\operatorname{Spec}(A)$, $A=\widehat{\mathbb{Z}}$,
\[
G=\mathrm{Gl}(2)_S,\quad
\rho=\begin{pmatrix}0&1\\-1&1\end{pmatrix},\quad
\sigma=\begin{pmatrix}0&-1\\1&0\end{pmatrix},
\]
and where we thus have a canonical homomorphism
\[
(2)\qquad
\widehat{\unsure{\mathcal{G}|\mathcal{C}}}_{0,3}
\longrightarrow
G(\widehat{\mathbb{Z}})^{\prime\pm1}
\simeq\unsure{\mathrm{Gl}|\mathrm{Sl}}(2,\widehat{\mathbb{Z}})^{\pm1}/\{\pm1\}
\]
with values in the subgroup of \unsure{elements|elements} of determinant
$\pm1$,
\margin{378}

% Printed archive footer (not autograph).
\begin{center}
378\\
\texttt{https://grothendieck.umontpellier.fr}\\
Université de Montpellier, UMR CNRS 5149, Case courrier 051, Place Eugène Bataillon\\
34095 Montpellier cedex 5 - France
\end{center}

\scanpage{380}

sending $\widehat{\unsure{\mathcal{G}|\mathcal{C}}}_{0,3}^{+}$ into
$SG(\widehat{\mathbb{Z}})'=\mathrm{Sl}(2,\widehat{\mathbb{Z}})'
=\mathrm{Sl}(2,\widehat{\mathbb{Z}})/\unsure{\{\pm1\}|\pm1}$,
we are going to construct a variant of the group $\mathcal{M}_{0,3}^{\sim}$,
\inserted{(cf. page 390)}, which maps canonically into the product (1),
which will lead us to define a group structure on
$S\mathcal{M}_{\rho,\sigma}$, so as to make the said morphism
multiplicative.

Let us first consider the ``special'' variant
$S\mathcal{M}_{0,3}^{\sim}$ of $\mathcal{M}_{0,3}^{\sim}$, an extension
of $\mathcal{M}_{0,3}^{\sim}$ by $\widehat{\mathbb{Z}}^3$,
deduced from an extension of
\[
\unsure{\mathcal{A}|\mathcal{H}}_{0,3}^{\sim}
=\mathcal{M}_{0,3}^{\sim}/\widehat{\Pi}_{0,3}
\subset\unsure{\widehat{\widehat{\mathcal{T}}}_{0,3}|\widehat{\mathcal{T}}_{0,3}},
\]
\deleted{\unsure{p}} and even from an extension of
$\unsure{\widehat{\widehat{\mathcal{T}}}_{0,3}|\widehat{\mathcal{T}}_{0,3}}$
(via $\mathcal{M}_{0,3}^{\sim}\longrightarrow
\unsure{\widehat{\widehat{\mathcal{T}}}_{0,3}|\widehat{\mathcal{T}}_{0,3}}$),
consisting of the
\[
(3)\qquad
(u,\unsure{f_0|g_0},\unsure{f_1|g_1},\unsure{f_\infty|g_\infty})
\qquad
\left(u\in\unsure{\widehat{\widehat{\mathcal{T}}}_{0,3}|\widehat{\mathcal{T}}_{0,3}},\quad
g_0,g_1,g_\infty\in\unsure{\widehat{\Pi}_{0,3}}\right)
\]
satisfying
\[
(4)\qquad
u(\ell_i)=\operatorname{int}\bigl(\unsure{f|g}_{\dot{u}(i)}^{-1}\bigr)
\bigl(\ell_{\dot{u}(i)}^{\mu}\bigr)
\]
where $\mu$ is the multiplier of $u$, and $\dot{u}$ the action of $u$
on \inserted{the set} $\{0,1,\infty\}$ of loop classes.
The product is defined by
\[
\begin{aligned}
(5)\qquad &(u',f'_0,f'_1,f'_\infty)
(\deleted{\unsure{u}},f_0,f_1,f_\infty)\\
&=(u'u,\ f'_0u'(f_0),\ f'_1u'(f_1),\ \unsure{f'_\infty u'(f_\infty)})
\end{aligned}
\]
restricting ourselves (\unsure{deflation!}) to the case
$u,u'\in\mathcal{M}_{0,3}^{!\sim}$ i.e. $\dot{u},\dot{u}'=1$.

\margin{379}

% Printed archive footer (not autograph).
\begin{center}
379\\
\texttt{https://grothendieck.umontpellier.fr}\\
Université de Montpellier, UMR CNRS 5149, Case courrier 051, Place Eugène Bataillon\\
34095 Montpellier cedex 5 - France
\end{center}

\scanpage{381}
\mspage{491}

We have a surjective homomorphism $S\mathcal{M}_{0,3}^{\sim}\longrightarrow
\mathcal{M}_{0,3}^{\sim}$, whose kernel consists of the
\[
(6)\qquad (1,f_0,f_1,f_\infty)\qquad f_i\in L_i
\]
\unsure{thus we have}
\[
(7)\qquad
\diagram{\begin{tikzcd}[ampersand replacement=\&,column sep=large]
1\arrow[r]\&\widehat{\mathbb{Z}}^3\arrow[r]\&
S\mathcal{M}_{0,3}^{\sim}\arrow[r]\&
\mathcal{M}_{0,3}^{\sim}\arrow[r]\&1
\end{tikzcd}}
\]

We are going to define an interesting subgroup
\[
(8)\qquad S_0\mathcal{M}_{0,3}^{\sim}\subset
\mathcal{M}_{0,3}^{\sim},
\]
now using the definition of $\mathcal{M}_{0,3}^{\sim}$,
namely \inserted{\unsure{via}} the property of the
$u\in\mathcal{M}_{0,3}^{\sim}$ of commuting with
$\rho$ as outer actions, so that we will have
\[
(9)\qquad u\rho u^{-1}=g\rho\qquad\text{i.e. }[u,\rho]=g,
\quad g\in\widehat{\Pi}_{0,3}
\]

This allows us, when we have already \uline{chosen}
$f_i\in\widehat{\Pi}_{0,3}$ satisfying
\[
(10)\qquad u(\ell_i)=\operatorname{int}(f_i^{-1})(\ell_i^{\mu})
\]
(supposing for simplicity $u\in\mathcal{M}_{0,3}^{!\sim}$),
to deduce from it the \deleted{\unsure{forms|formulas}}
$f_{i+1},f_{i+2}$ via the formula
\[
\begin{aligned}
(11)\qquad f_{i+1}^{-1}&=g\rho(f_i^{-1})\,,\qquad
\inserted{\left[\begin{aligned}
\text{and }f_{i+2}^{-1}&=g\rho(f_{i+1}^{-1})\\
&=g\rho(g)\rho^2(f_i^{-1})
\end{aligned}\right]}\\[-1ex]
&\inserted{[\text{where }g=[u,\rho]\in\widehat{\Pi}_{0,3})}
\end{aligned}
\]
(we have set $\dot{\rho}(i)=i+1$ by abuse of language,
i.e. $0+1=1,\ 1+1=\infty,\ \infty+1=0$!) -- we
then \unsure{recover|will recover}
\[
f_i^{-1}=\unsure{f|g}_{(i+2)+1}^{-1}
=g\rho(f_{i+2}^{-1})
=\underbrace{g\rho(g)\rho^2(g)}_{1}\,
\underbrace{\rho^3(f_i^{-1})}_{f_i^{-1}}.
\]
\margin{380}

% Printed archive footer (not autograph).
\begin{center}
380\\
\texttt{https://grothendieck.umontpellier.fr}\\
Université de Montpellier, UMR CNRS 5149, Case courrier 051, Place Eugène Bataillon\\
34095 Montpellier cedex 5 - France
\end{center}

\scanpage{382}
\mspage{492}

The $(u,f_0,f_1,f_\infty)
\in S\mathcal{M}_{0,3}^{!\sim}$ which satisfy (11) for
$\forall\ i\in\{0,1,\infty\}$ (it suffices that it satisfy \unsure{it|them}
for two indices, for this to be satisfied for all
three) form a \uline{subgroup} of
$S\mathcal{M}_{0,3}^{!\sim}$, which I denote
$S_0\mathcal{M}_{0,3}^{!\sim}$ -- it also identifies with the group
of pairs $(u,f_0)$ $(u\in\mathcal{M},\ f_0\in\widehat{\Pi}_{0,3})$
satisfying
\[
(12)\qquad u(\ell_0)=\operatorname{int}(f_0^{-1})(\ell_0^{\mu})
\qquad(\inserted{\text{where }}\mu=\chi(u)).
\]

We have an extension structure
\[
(13)\qquad
\diagram{\begin{tikzcd}[ampersand replacement=\&,column sep=large]
1\arrow[r]\&\widehat{\mathbb{Z}}\arrow[r,no head]\&
S_0\mathcal{M}_{0,3}^{!\sim}\arrow[r]\&
\mathcal{M}_{0,3}^{!\sim}\arrow[r,no head]\&1
\end{tikzcd}},
\]
and a homomorphism of extensions
\[
(14)\qquad\left\{
\diagram{\begin{tikzcd}[ampersand replacement=\&,column sep=large,row sep=large]
1\arrow[r,no head]\&\widehat{\mathbb{Z}}\arrow[r,no head]\arrow[d,"\text{diagonal}"]\&
S_0\mathcal{M}_{0,3}^{!\sim}\arrow[r,no head]\arrow[d]\&
\mathcal{M}_{0,3}^{!\sim}\arrow[r,no head]\arrow[d]\arrow[d,no head,shift right=0.5ex]\&1\\
1\arrow[r,no head]\&\widehat{\mathbb{Z}}^3\arrow[r,no head]\&
S\mathcal{M}_{0,3}^{!\sim}\arrow[r,no head]\&
\mathcal{M}_{0,3}^{!\sim}\arrow[r,no head]\&1
\end{tikzcd}}\right.
\]
(where $S\mathcal{M}_{0,3}^{!\sim}$ is \inserted{defined as}
the \unsure{reverse|inverse} image of
$\mathcal{M}_{0,3}^{!\sim}\unsure{\subset|\in}\mathcal{M}_{0,3}^{\sim}$
in $S\mathcal{M}_{0,3}^{\sim}$). We can therefore say that
the extension $S\mathcal{M}_{0,3}^{!\sim}$ (of
$\mathcal{M}_{0,3}^{!\sim}$ by $\widehat{\mathbb{Z}}^3$) is
deduced, by the diagonal homomorphism
\[
(15)\qquad
\diagram{\begin{tikzcd}[ampersand replacement=\&,column sep=large,row sep=small]
\widehat{\mathbb{Z}}\arrow[r]\&\widehat{\mathbb{Z}}^3\\
\underline{n}\arrow[r,mapsto]\&
(\underline{n}_0=\underline{n},\ \underline{n}_1=\underline{n},\
\underline{n}_\infty=\underline{n})
\end{tikzcd}},
\]
from a ``finer'' extension $S_0\mathcal{M}_{0,3}^{!\sim}$ of
\margin{381}

% Printed archive footer (not autograph).
\begin{center}
381\\
\texttt{https://grothendieck.umontpellier.fr}\\
Université de Montpellier, UMR CNRS 5149, Case courrier 051, Place Eugène Bataillon\\
34095 Montpellier cedex 5 - France
\end{center}

\scanpage{383}
$\mathcal M_{0,3}^{!\sim}$ by $\widehat{\mathbb Z}$.

The calculations of \S 46 \uline{V} (page 407 \unsure{and following}) make it possible on the other hand to make explicit the pairs $(\unsure{u|\alpha},f_0=f)$ in terms of systems
\[
(16)\qquad
(\alpha,\beta,f)\quad\left|\quad
\begin{gathered}
\alpha\in\widehat{\widetilde{\Pi}}_{0,3}
 =\widehat{\Pi}_{0,3}\cdot\{1,\tau_\infty\}
 =\text{\unsure{inverse image|inverse image} of }\{1,\tau\}\\
\text{under }\widehat{\unsure{\mathcal G|\mathcal C}}_{0,3}\longrightarrow\mathcal T_{0,3}\\
\beta,f\in\widehat{\Pi}_{0,3}
\end{gathered}\right.
\]
\margin{to $(u,f)$ we associate $(\unsure{fu},f)$ by putting $\unsure{fu}\in\mathcal M_{0,3}^{!\sim}[0]$ in the form $u_{\alpha,\beta}$}
satisfying relations (not involving $f$)
\[
(17)\qquad
\alpha\rho(\alpha)^{-1}
 =\beta\sigma(\beta)^{-1}
 \underbrace{\varepsilon_1(\mu-1)}_{=\unsure{\ell_1|h}^{\nu}\ \text{if }\mu-1=2\nu}
\]
where $\mu\in\widehat{\mathbb Z}^{*}$ is defined \deleted{\unsure{without any}} \unsure{ambiguity|unambiguously} by $\alpha$ --- and even \unsure{from knowledge alone} \deleted{of the image} of $\xi=\underline\alpha\rho(\underline\alpha)^{-1}$, where $\underline\alpha$ is the image of $\alpha$ in $\mathrm{Gl}(2,\widehat{\mathbb Z})^{\pm1}$.

The case $f=1$ corresponds to \inserted{a} subgroup of $S_0\mathcal M_{0,3}^{!\sim}$, which is identified with the normalizer of $L_0$ in $\mathcal M_{0,3}^{!\sim}$, which we had denoted $\mathcal M_{0,3}^{!\sim}[0]$ (\S 46 \uline{V} $\ldots$) :
\[
\begin{gathered}
(18)\qquad
\underset{\substack{\Vert\ \mathrm{dfn}\\
\operatorname{Norm}_{\mathcal M_{0,3}^{!\sim}}(L_0)}}{\mathcal M_{0,3}^{!\sim}[0]}
\simeq
\left\{(\alpha,\beta,f)\,\middle|\,
\begin{array}{l}f=1,\ \alpha,\beta\\\text{satisfying }(17)\end{array}\right\}
\subset S_0\mathcal M_{0,3}^{!\sim} .
\end{gathered}
\]
If we denote by $\unsure{u_{\xi,\eta}|u_{\alpha,\beta}}$ \deleted{\unsure{or $u_\xi$}} \deleted{the automorphism} \unsure{or simply} $u_\xi$ the automorphism defined in terms of
\[
(19)\qquad
\xi=\alpha\rho(\alpha)^{-1},\qquad
\eta=\beta\sigma(\beta)^{-1}
\]
\margin{382}

% Printed archive footer (not autograph).
\begin{center}
382\\
\texttt{https://grothendieck.umontpellier.fr}\\
Université de Montpellier, UMR CNRS 5149, Case courrier 051, Place Eugène Bataillon\\
34095 Montpellier cedex 5 - France
\end{center}

\scanpage{384}
\mspage{493}
by
\[
(20)\qquad u(\rho)=\xi\rho,\qquad u(\sigma)=\eta\sigma
\]
(N.B. $\eta$ is defined in terms of $\xi$ thanks to (17), where $\mu=\mu(\xi)$ is determined by $\xi$), we \unsure{will have} by construction (loc. cit.)
\[
(21)\qquad u_{\alpha,\beta,f}=\operatorname{int}(f^{-1})u_{\alpha,\beta}
\]
and to the triple $(\alpha,\beta,f)$ corresponds the element $(u_{\alpha,\beta,f},f)$ of $S_0\mathcal M_{0,3}^{!\sim}$ :
\[
(22)\qquad (\alpha,\beta,f)\longmapsto (u_{\alpha,\beta},f)\in S_0\mathcal M_{0,3}^{!\sim} .
\]
Having thus made the elements of $S_0\mathcal M_{0,3}^{!\sim}$ explicit, we define a map
\[
(23)\qquad
\diagram{\begin{tikzcd}[ampersand replacement=\&,column sep=large,row sep=large]
S_0\mathcal M_{0,3}^{!\sim}\arrow[r] \&
S\mathcal M_{\rho,\sigma}(\widehat{\mathbb Z})'
=S\mathcal G_{\rho,\sigma}^{*}(\widehat{\mathbb Z})'
\times\mathrm{Sl}(2,\widehat{\mathbb Z})\unsure{'}\\
(\alpha,\beta,f)\arrow[r,mapsto] \&
\left(\underline\alpha,\underline\beta,
\mu=\mu(\underline\xi=\underline\alpha\rho(\underline\alpha)^{-1}),
\underset{\substack{\uparrow\\\inserted{\unsure{\varepsilon}}}}{\varepsilon_2(\mu)},\underline f\right)
\end{tikzcd}}
\]
where the underlined $\underline\alpha,\underline\beta,\underline f$ \unsure{\&} denote the images of $\alpha,\beta,f$ in $\mathrm{Sl}(2,\widehat{\mathbb Z})'$, and the \unsure{``accent''} $'$ in $S\mathcal M_{\rho,\sigma}(\widehat{\mathbb Z})'$ denotes the set of \deleted{pairs} \inserted{systems} $(\underline\alpha,\underline\beta,\mu,\varepsilon,\underline f)$
\deleted{\unsure{defined of elements of}}
\[
\in\mathrm{Gl}(2,\widehat{\mathbb Z})^{\prime\unsure{\pm1}}
\times\mathrm{Sl}(2,\widehat{\mathbb Z})'
\times\widehat{\mathbb Z}^{*}\times\{\pm1\}
\times\inserted{\mathrm{Sl}(2,\widehat{\mathbb Z})'}
\]
satisfying \inserted{\uline{in $\mathrm{Sl}(2,\widehat{\mathbb Z})$} (\uline{not in $\unsure{\mathrm{Sl}}(2,\widehat{\mathbb Z})'$ !}),} the loop relation \deleted{\unsure{taken|by} modulo passage to}
\[
(24)\qquad
\underline\alpha\rho(\underline\alpha)^{-1}
 =\varepsilon\underline\beta\sigma(\underline\beta)^{-1}\varepsilon_1(\mu-1)
\]
(which does not depend on the choices of \uline{signs} $\ldots$)
\margin{383}

% Printed archive footer (not autograph).
\begin{center}
383\\
\texttt{https://grothendieck.umontpellier.fr}\\
Université de Montpellier, UMR CNRS 5149, Case courrier 051, Place Eugène Bataillon\\
34095 Montpellier cedex 5 - France
\end{center}

\scanpage{385}
We have ``done what was needed'' so that the restriction of (23) to $\deleted{\unsure{S_0}}\mathcal M_{0,3}^{!\sim}[0]$ ((18), identified with the $(\alpha,\beta,f)$ for which $f=1$) is a homomorphism \inserted{into the group $S\mathcal G_{\rho,\sigma}$} --- by \unsure{defining} precisely in a clever way \inserted{ad hoc} the composition law in $S\mathcal G_{\rho,\sigma}^{*}$.
I am going to recheck it all the same to \unsure{reassure myself}, by considering $u=u_{\alpha,\beta}$, $u'=u_{\alpha',\beta'}$, $u''=uu'$ in the form $u_{\alpha'',\beta''}$, and recalling the ``cocycle relation'' which gives, up to one of the two \unsure{signs}, $(\alpha'',\beta'')$ in terms of $(\alpha',\beta'),(\alpha,\beta)$ (\S 46 \uline{II} $5^\circ$))
\[
(25)\qquad
\alpha''=u'_{\unsure{\infty|\omega}}(\alpha)\alpha',\qquad
\beta''=u'(\beta)\beta',\qquad
\inserted{\text{in addition to }\mu''=\mu\mu'}
\]
Let \unsure{then|therefore} $\underline\alpha,\underline\beta,\underline\alpha',\underline\beta'$ be \deleted{images} \inserted{\unsure{elements}} of $\mathrm{Gl}(2,\widehat{\mathbb Z})^{\pm1}$ (determined up to \uline{sign}) which lift the images of $\alpha,\alpha',\beta,\beta'$ in $\mathrm{Gl}(2,\widehat{\mathbb Z})'$, we will therefore have $(\underline\alpha,\underline\beta,\mu,\varepsilon_2(\mu))$ and $(\underline\alpha',\underline\beta',\mu',\varepsilon_2(\mu'))$ in $S\mathcal G_{\rho,\sigma}^{*}(\widehat{\mathbb Z})$, and applying $\Theta$ we find $\underline u,\underline u'\in\mathrm{Gl}(2,\widehat{\mathbb Z})$ (in fact, $\in B'_0(\widehat{\mathbb Z})$) such that
\[
(26)\qquad
\left\{\begin{aligned}
\underline u'(\rho)&=\operatorname{int}(\underline\alpha')\rho
&\underline u'(\sigma)&=\varepsilon_2(\mu')\operatorname{int}(\underline\beta')(\sigma)\\
\det\underline u'&=\mu'
\end{aligned}\right.
\]
(and relations \unsure{item|idem} for $u$)
and we will then have, for every $x\in\deleted{\unsure{\Pi}}\widehat{\unsure{\mathcal G|\mathcal C}}_{0,3}$
\[
(27)\qquad \underline{u'(x)}=\underline u'(\underline x)
\]
(cf. $\ldots$) Thus, for \deleted{some} suitable liftings $\underline\alpha'',\underline\beta^{\unsure{''|'}}$ of the images of $\alpha^{\unsure{''|'}},\beta^{\unsure{''|'}}$ in $\mathrm{Gl}(2,\widehat{\mathbb Z})^{\pm1}$, we will have
\[
\begin{gathered}
(28)\qquad
\underline\alpha''=\underline u'(\underline\alpha)\underline\alpha',\qquad
\underline\beta''=\underline u'(\underline\beta)\underline\beta',\\
\text{in addition to }\mu''=\mu\mu'\quad\text{hence}\quad
\varepsilon_2''=\varepsilon_2\varepsilon_2'.
\end{gathered}
\]
\margin{384}

% Printed archive footer (not autograph).
\begin{center}
384\\
\texttt{https://grothendieck.umontpellier.fr}\\
Université de Montpellier, UMR CNRS 5149, Case courrier 051, Place Eugène Bataillon\\
34095 Montpellier cedex 5 - France
\end{center}

\scanpage{386}
\mspage{494}
and it remains to check that the element
\[
\begin{gathered}
g''=(\underline\alpha'',\underline\beta'',\mu''=\mu\mu',\varepsilon_2''=\varepsilon_2'\varepsilon_2)
\in\deleted{\unsure{\Gamma}(\mathrm{Gl}(2,\widehat{\mathbb Z})\times\mathrm{Sl}(2,\widehat{\mathbb Z})\times}\\
\inserted{\Gamma\bigl(\widehat{\mathbb Z};\mathrm{Gl}(2)\times\mathrm{Sl}(2)\times\unsure{\mathbf E}'\times\mu_2\bigr)}
\end{gathered}
\]
is indeed the product of the elements
\[
g'=(\underline\alpha',\underline\beta',\mu',\varepsilon')\quad\text{and}\quad
g=(\underline\alpha,\underline\beta,\mu,\varepsilon).
\]
That it is OK for $\mu''$ and $\varepsilon''$ is trivial, it remains to check that in $\mathcal G_{\rho,\sigma}^{*}$, the product is indeed (as regards the components in $\alpha$ and in $\beta$) given by formulas (28) --- which \unsure{will be} valid, of course, over any base, and under the general conditions of \uline{V} (or \inserted{even} of \uline{III}, if we are interested in $\mathcal G_{\rho,\sigma}$, instead of $\mathcal G_{\rho,\sigma}^{*}$). We will have
\[
(29)\qquad
\left\{\begin{aligned}
\underline\alpha&=\underline u\,\underline a^{-1},&
\underline\beta&=\underline u\,\underline b^{-1}\\
\underline\alpha'&=\underline u'\,\underline a'^{-1},&
\underline\beta'&=\underline u'\,\underline b'^{-1}
\end{aligned}\right.
\qquad
\begin{gathered}
\underline\alpha,\underline\alpha'\in T_\rho(\widehat{\mathbb Z})\\
\underline\beta,\underline\beta'\in N'_\sigma(\widehat{\mathbb Z})
\end{gathered}
\]
and by definition of the products in $\mathcal G_{\rho,\sigma}^{*}$, the product of \deleted{\unsure{computed at}} $g'g$ will therefore have components $\underline\alpha_0''$ and $\underline\beta_0''$
\[
\underline\alpha_0''=(\underline u'\underline u)(\underline a'\underline a)^{-1},\qquad
\underline\beta_0''=(\underline u'\underline u)(\underline b'\underline b)^{-1}
\]
and it remains to prove
\[
\underline\alpha_0''=\underline\alpha'',\qquad
\underline\beta_0''=\underline\beta'',\quad\text{i.e. the relations}
\]
\[
(30)\qquad
\underbrace{\underline u'(\underline\alpha)\cdot\underline\alpha'}
 =(\underline u'\underline u)(\underline a'\underline a)^{-1}
\qquad
\underline u'(\underline\beta)\underline\beta'
 =(\underline u'\underline u)(\underline b'\underline b)^{-1}
\]
where $\deleted{\unsure{\underline u}}\ \underline\alpha,\underline\beta,\underline\alpha',\underline\beta'$ are given in terms of \deleted{$\unsure{\underline\alpha,\underline\beta,\underline u,\underline a}$} $\underline u,\underline a,\underline b,\underline u',\underline a',\underline b'$ by (29). We therefore have
\[
\begin{aligned}
\underline u'(\underline\alpha)\underline\alpha'
 &=\underbrace{\underline u'(\underline u\,\underline a^{-1})\underline u'^{-1}}\underline u'\underline a'^{-1}\\
 &=(\underline u'\underline u)(\underline a^{-1}\underline a'^{-1})
 =(\underline u'\underline u)(\underline a'\underline a)^{-1},
 \qquad\deleted{\text{\unsure{OK}}}
\end{aligned}
\]
\deleted{$\unsure{\underline u'(\underline\beta)\underline\beta'}$} and likewise for the second relation 30, which proves our assertion.
\margin{$\unsure{\checkmark|/}$}
\margin{385}

% Printed archive footer (not autograph).
\begin{center}
385\\
\texttt{https://grothendieck.umontpellier.fr}\\
Université de Montpellier, UMR CNRS 5149, Case courrier 051, Place Eugène Bataillon\\
34095 Montpellier cedex 5 - France
\end{center}

\scanpage{387}
Let us note on this occasion that the systems
\[
(31)\qquad
(\alpha,\beta,\unsure{\mu}=1,\varepsilon=1)
\quad\text{with}\quad
\alpha\in\Gamma\deleted{\unsure{\{\pm1\}}}\inserted{\mathbb G_m},\qquad
\beta\in\Gamma\deleted{\unsure{\{\pm1\}}}\inserted{\mathbb G_m}
\]
are in $\unsure{\Gamma|S}\mathcal G_{\rho,\sigma}^{\unsure{\circ|0}}$.
\deleted{and even in $S\mathcal G_{\rho,\sigma}$ if $\alpha,\beta\in\Gamma\mu_2$ ---}
[which are indeed stable under multiplication, \deleted{\unsure{for}} $(\alpha,\beta,\mu,\varepsilon)$, of $\alpha$ and $\beta$ each by a section of $\deleted{\unsure{\mu_2}}\inserted{\mathbb G_m}$, which does not change $\widetilde\alpha,\widetilde\beta\ \ldots$]
they correspond, by the isomorphism $\mathcal G_{\rho,\sigma}\simeq B'_0\times T_\rho\times T_\sigma$, to the systems $(\underline u,\underline a,\underline b)$ with $\underline u=1$, $\underline a\in\Gamma\mathbb G_m$, $\underline b\in\Gamma\mathbb G_m$) and form a \uline{subgroup} (with ordinary multiplication \deleted{\unsure{on the|on the}} \deleted{in $\alpha,\beta$} \inserted{in $\alpha\ \beta$} as group law) --- if $\alpha,\beta$ are sections of $\mu_2$, we find a \unsure{subgroup} \deleted{trivial} ``\unsure{by nature|in action} trivial'' \inserted{$\mu_2\times\mu_2$} of $S\mathcal G_{\rho,\sigma}$.
The action by translation of this subgroup $\simeq\mathbb G_m\times\mathbb G_m$ on $\mathcal G_{\rho,\sigma}^{*}$ is given by multiplication \deleted{in} \inserted{of the} \deleted{the} components in question. This shows that the fact of working with $\underline\alpha,\underline\beta$ defined only ``up to \uline{signs}'' simply amounts to working in the quotient groups of \uline{$\mathcal G_{\rho,\sigma}^{*}(\widehat{\mathbb Z})$} (or of $S\mathcal G_{\rho,\sigma}^{*}(\widehat{\mathbb Z})$) by the subgroup
\[
\inserted{\boxed{\begin{gathered}\Sigma_0\ .\\ \simeq\mathbb Z/2\mathbb Z\times\mathbb Z/2\mathbb Z\end{gathered}}}
\]
\deleted{generated} of the elements $(\alpha,\beta,1,1)$ with $\alpha,\beta\in\{\pm1\}$. This
\margin{386}

% Printed archive footer (not autograph).
\begin{center}
386\\
\texttt{https://grothendieck.umontpellier.fr}\\
Université de Montpellier, UMR CNRS 5149, Case courrier 051, Place Eugène Bataillon\\
34095 Montpellier cedex 5 - France
\end{center}

\scanpage{388}
\mspage{495}

makes precise at the same time the meaning of the expression:
we have a \unsure{homomorphism} of groups \ldots .

Let us now return to the map (23) with arbitrary $f$, and make precise the
composition law of $S_0\mathcal M_{0,3}^{!\sim}$ in terms of
$(\alpha,\beta,f)$ i.e. in terms of the $u_{\alpha,\beta,f}$.
Let us consider the subgroup formed of the $(1,1,f)$,
\deleted{or again} \deleted{of the} which corresponds to the \inserted{pairs}
$(u,f)$ with $u=\operatorname{int}(f^{-1})$, i.e. to the pairs
$(\operatorname{int}(f^{-1}),f)$
--- they form a \unsure{\textquotedblleft}group isomorphic to $\widehat\Pi_{0,3}$,
and we have a $\frac12$ direct product structure:
\[
(31)\qquad
S_0\mathcal M_{0,3}^{!\sim}
\simeq \mathcal M_{0,3}^{!\sim}[0]\cdot\widehat\Pi_{0,3}
\]
where
\[
(32)\qquad
\left\{
\diagram{\begin{tikzcd}[ampersand replacement=\&,column sep=large,row sep=small]
\widehat\Pi_{0,3} \arrow[r] \&
S_0\mathcal M_{0,3}^{!\sim} \& {}\\
f \arrow[r,mapsto] \&
(\deleted{\unsure{f,}}\operatorname{int}(f),f)
\arrow[r,leftrightarrow] \& (1,1,f^{-1}),
\end{tikzcd}}
\right.
\]
the operation of $\mathcal M_{0,3}^{!\sim}[0]$ on $\widehat\Pi_{0,3}$
\deleted{\unsure{being}} the tautological operation
(deduced from the inclusion
$\mathcal M_{0,3}^{!\sim}[0]\subset\mathcal M_{0,3}^{!\sim}
\subset\underline{\operatorname{Aut}}_{\mathrm{lac}}
(\widehat\Pi_{0,3})$).

The \inserted{invariant} subgroup $\widehat{\mathbb Z}$ of
$S_0\mathcal M_{0,3}^{!\sim}$ is given by the pairs
$\inserted{(u,f)=}(1,\ell_0^n)$ $(n\in\widehat{\mathbb Z})$
which we also put in the form $(\alpha,\beta,f)$
where $\alpha,\beta$ are such that
$u_{\alpha,\beta}=fu=\operatorname{int}(\ell_0^n)$,
hence $\alpha=\beta=\ell_0^n$, i.e.
\[
(33)\qquad
\widehat{\mathbb Z}\longrightarrow S_0\mathcal M_{0,3}^{!\sim}
\quad\text{is identified with}\quad
n\longmapsto
(\deleted{\unsure{\alpha=}}\ell_0^n,\beta=\ell_0^n,f=\ell_0^n)
\]
\margin{387}

% Printed archive footer (not autograph).
\begin{center}
387\\
\texttt{https://grothendieck.umontpellier.fr}\\
Université de Montpellier, UMR CNRS 5149, Case courrier 051, Place Eugène Bataillon\\
34095 Montpellier cedex 5 - France
\end{center}

\scanpage{389}

We \unsure{will have}
\[
\begin{aligned}
(f'^{-1}u_{\alpha',\beta'})(f^{-1}u_{\alpha,\beta})
&=\deleted{\unsure{f'^{-1}f^{-1}(f u_{\alpha',\beta'}f^{-1})}}\\
&\inserted{\underbrace{f'^{-1}\unsure{u_{\alpha',\beta'}}(f^{-1})}
_{\unsure{f''^{-1}}}
(u_{\alpha',\beta'}u_{\alpha,\beta})}
\end{aligned}
\]
\unsure{i.e.} the multiplication in $S_0\mathcal M_{0,3}^{!\sim}$
is given, in terms of the components
$(\alpha',\beta',f')(\alpha,\beta,f)$,
\unsure{for|by} \deleted{\unsure{g''=}}$(\alpha'',\beta'',f'')$, where
\[
(34)\quad
\left\{
\begin{aligned}
u_{\alpha'',\beta''}
&=u_{\alpha',\beta'}u_{\alpha,\beta}
\quad\text{i.e.}\quad
\left\{
\begin{aligned}
\alpha''&=u_{\alpha',\beta'}(\alpha)\alpha'\\
\beta''&=u_{\alpha',\beta'}(\beta)\beta'
\end{aligned}
\right.,
\quad\text{and}\\
f''&=u_{\alpha',\beta'}(f)f'
\end{aligned}
\right.
\]

Returning now to the definition of $S\mathcal M_{\rho\sigma}$
(1), we are going to introduce there the \unsure{corresponding} multiplication,
which amounts to defining $S\mathcal M_{\rho\sigma}$ as
a $\frac12$ direct product
\[
\begin{aligned}
(35)\qquad S\mathcal M_{\rho,\sigma}
&\overset{\mathrm{dfn}}{=}
S\mathcal G_{\rho,\sigma}^{*}\cdot SG\\
&\qquad(\text{product }\tfrac12\text{ direct})
\end{aligned}
\]
\uline{by the operation of} $S\mathcal G_{\rho,\sigma}^{*}$
\uline{on} $G$
\uline{via the homomorphism} $\Theta$ (page~\unsure{489'|488'}),
and \unsure{by} \unsure{defining|will define}
\unsure{the|the} \deleted{map}\inserted{morphism}
\[
(36)\qquad
\left|
\diagram{\begin{tikzcd}[ampersand replacement=\&,column sep=large,row sep=small]
S\mathcal M_{\rho,\sigma}\arrow[r] \&
S\mathcal G_{\rho,\sigma}^{*}\times_S G\\
\deleted{\unsure{(\alpha,\beta,\mu,\varepsilon)|(\alpha,\beta)}}
\inserted{U}\cdot f\arrow[r,mapsto] \& (U,f^{-1})
\end{tikzcd}}
\right.
\]
\margin{Watch out for the passage from $f\underset{\leftarrow}{\ \unsure{\text{to}}\ }f^{-1}$!}
where $U=(\alpha,\beta,\mu,\varepsilon)$ is in
$\Gamma S\mathcal G_{\rho,\sigma}$.
\margin{388}

% Printed archive footer (not autograph).
\begin{center}
388\\
\texttt{https://grothendieck.umontpellier.fr}\\
Université de Montpellier, UMR CNRS 5149, Case courrier 051, Place Eugène Bataillon\\
34095 Montpellier cedex 5 - France
\end{center}

\scanpage{390}
\mspage{496}

We have thus done what was needed, in the case
$S=\operatorname{Spec}(\widehat{\mathbb Z})$ etc, to have a
\uline{homomorphism of groups}
\[
(37)\qquad
\diagram{\begin{tikzcd}[ampersand replacement=\&,column sep=large]
S_0\mathcal M_{0,3}^{!\sim}\arrow[r] \&
S\mathcal M_{\rho,\sigma}(\widehat{\mathbb Z})/
\unsure{\widetilde{\Sigma}}\deleted{\unsure{(\widehat{\mathbb Z})}}
\end{tikzcd}}
\]
where $\Sigma_0\simeq\mathbb Z_{/2}\times\mathbb Z_{/2}$
is the \unsure{subgroup} \deleted{\unsure{\textquotedblleft powerless\textquotedblright}}
of $S\mathcal G_{\rho,\sigma}^{*}$ introduced on page~494',
which operates trivially on $G$ hence is invariant in
$S\mathcal M_{\rho,\sigma}$, hence in
$S\mathcal M_{\rho,\sigma}(\widehat{\mathbb Z})$, and where
\[
(38)\quad
\widetilde{\Sigma}=\Sigma_0\times
\underbrace{\{\pm1\}}_{
\substack{
\text{this factor identified with the center of }SG(\unsure{\widehat{\mathbb Z}})\\
=\operatorname{Sl}(2,\widehat{\mathbb Z}),\
\text{contained in }\operatorname{Gl}(2,\widehat{\mathbb Z})\\
=G(\widehat{\mathbb Z})
}}
\simeq
\mathbb Z/2\mathbb Z\times\mathbb Z/2\mathbb Z\times\mathbb Z/2\mathbb Z
\]
is the corresponding normal group in the $\frac12$ direct product
--- the fact of passing to the quotient by $\Sigma$ expresses
the three \unsure{ambiguities|ambiguities} of \uline{signs},
on $\alpha$, on $\beta$ and on $f$.

Let us now consider the \unsure{homomorphism}
\[
(39)\quad
\left|
\diagram{\begin{tikzcd}[ampersand replacement=\&,column sep=large,row sep=small]
\mathbb G_a\arrow[r] \&
\unsure{S_0|S}\mathcal M_{\rho,\sigma}^{\deleted{\unsure{!\sim}}}\\
\lambda\arrow[r,mapsto] \&
(\alpha\unsure{=}\varepsilon_0(\unsure{\lambda|-\lambda}),
\beta\unsure{=}\varepsilon_0(-\lambda),
\mu=1,\varepsilon=1,f=\varepsilon_0(\lambda))
\end{tikzcd}}
\right.
\]
it is a homomorphism of groups, and even makes
$\mathbb G_a$ a \uline{normal} subgroup.
To \inserted{verify} this the most convenient thing is to identify
$\unsure{S_0|S}\mathcal M_{\rho,\sigma}^{!\sim}$
with the $\frac12$-direct product (35),
\margin{389}

% Printed archive footer (not autograph).
\begin{center}
389\\
\texttt{https://grothendieck.umontpellier.fr}\\
Université de Montpellier, UMR CNRS 5149, Case courrier 051, Place Eugène Bataillon\\
34095 Montpellier cedex 5 - France
\end{center}

\scanpage{391}

where we identify $S\mathcal G_{\rho,\sigma}^{*}$ with a subgroup
of $B'_0\times T_\rho\times N'_\sigma$, operating on $G$
by inner automorphisms \uline{via} $B'_0$.

The preceding morphism thus corresponds to
\[
(40)\qquad
\diagram{\begin{tikzcd}[ampersand replacement=\&,column sep=small,row sep=small]
\lambda\arrow[r,mapsto,"i"]\arrow[d,phantom,"\in"{rotate=90}] \&
{(\lambda,1,1,} \&
\lambda^{-1}\arrow[d,phantom,"\in"{rotate=90}] \&
{)=\deleted{\lambda\cdot\unsure{j}(\lambda)}}\\
\Gamma B'_0 \& {} \& \Gamma G \& {}
\end{tikzcd}}
\]
\margin{\deleted{\diagram{\begin{tikzcd}[ampersand replacement=\&,column sep=small]
\unsure{j}:SG\arrow[r] \& S\mathcal M_{\rho,\sigma}
\end{tikzcd}} is} the canonical homomorphism of groups}

\margin{\deleted{\unsure{where $U_0$}}}
\deleted{Let on the other hand
$U=(\underbrace{\underline u,\underline{\unsure{a|\rho}},
\underline{\unsure{b|\beta}}}_{U_0},\unsure{\varepsilon}f)$
\inserted{$=U_0\cdot f$},
be an arbitrary element of the $\frac12$ direct product,
and let us compute $U i(\lambda)U^{-1}$ ---
\[
U i(\lambda)U^{-1}=U_0 f i(\lambda)f^{-1}U_0^{-1}=
\]
}

We identify $S\mathcal G_{\rho,\sigma}^{*}$ with a subgroup of
$\mathcal G_{\rho,\sigma}^{*}$, and likewise
the $\frac12$ direct product $S\mathcal M_{\rho,\sigma}^{\unsure{*}}$
with a subgroup of the $\frac12$ direct product
\[
(41)\qquad
\deleted{\unsure{S\mathcal M_{\rho,\sigma}^{*}}}
\inserted{\widetilde S\mathcal M_{\rho,\sigma}}
=\mathcal G_{\rho,\sigma}^{*}\cdot G
\simeq(B'_0\times T_\rho\times N'_\sigma)\cdot G,
\]
and we are going to see that the image of the \unsure{monomorphism} $i$ (40) in
$\mathcal G_{\rho,\sigma}^{*}\cdot G$ is even
normal in this larger group.
So as not to get lost in the notation,
let us denote by
\[
(42)\qquad i_B,\ i_\rho,\ i_\sigma,\ i_G
\]
the morphisms \unsure{of inclusions|of inclusion}
\[
\left|
\diagram{\begin{tikzcd}[ampersand replacement=\&,column sep=large,row sep=small,nodes in empty cells]
B'_0\arrow[r] \& \deleted{\unsure{\mathcal G_{\rho,\sigma}^{*}}}\\
T_\rho\arrow[r] \& {}\\
\unsure{N'_\sigma|N''_\sigma}\arrow[r] \& {}\\
G\arrow[r] \& {}
\end{tikzcd}}
\right\}\ 
\unsure{i'}
\begin{gathered}
\inserted{\widetilde S\mathcal M_{\rho,\sigma}}\\
\deleted{\unsure{\mathcal G_{\rho,\sigma}^{*}|\mathcal M_{\rho,\sigma}}}\cdot G
\end{gathered}
\]
\margin{390}

% Printed archive footer (not autograph).
\begin{center}
390\\
\texttt{https://grothendieck.umontpellier.fr}\\
Université de Montpellier, UMR CNRS 5149, Case courrier 051, Place Eugène Bataillon\\
34095 Montpellier cedex 5 - France
\end{center}

\scanpage{392}
\mspage{497}

it being understood, of course, that we let
$\mathcal G_{\rho,\sigma}^{*}$ operate on $G$ by inner automorphisms via the
factor $B'_0$. We must check that
\[
i_{\deleted{\unsure{B}}}(L_0)
=\{i_B(\lambda)i_G(\lambda^{-1})\mid\lambda\in\Gamma L_0\}
\]
is invariant under
$\deleted{\unsure{\Gamma T_\rho},\unsure{\Gamma T_\sigma}\unsure{i_T}}$
\inserted{\unsure{the}}
$i_\rho(x)$, $i_\sigma(x)$, $i_B(x)$ and $i_G(x)$.
--- For the first two, this is trivial --- with trivial action on $L_0$.

If $x\in\Gamma B'_0$, we will have
\[
\begin{aligned}
i_B(x)(i(\lambda))
&=i_B(x)(i_B(\lambda)i_G(\lambda^{-1}))=\\
&=\underbrace{i_B(x)(i_B(\lambda))}
_{i_B(\deleted{\unsure{\lambda|x}}x(\lambda))}
\cdot
\underbrace{i_B(x)(i_G(\lambda^{-1}))}
_{\substack{
i_G(x)(i_G(\lambda^{-1}))\\
=i_G(x(\lambda)^{-1}]
}}\\
&=i(x(\lambda))
\end{aligned}
\]
\margin{We \unsure{see} at once that $i(\lambda)$
commutes with $G$, hence is invariant under $G$ !}
which shows that $\deleted{\unsure{\mathbb G_a|L_0}}$
\[
\diagram{\begin{tikzcd}[ampersand replacement=\&,column sep=large]
L_0\arrow[r,"i"] \&
\mathcal G_{\rho,\sigma}^{*}\cdot G
=\widetilde S\mathcal G_{\rho,\sigma}^{*}
\end{tikzcd}}
\]
commutes with the natural action of $B'_0$ on both sides
(on the second via $\unsure{\pi_B}$,
\[
\diagram{\begin{tikzcd}[ampersand replacement=\&,column sep=large]
i_B:B'_0\arrow[r] \&
\widetilde S\underset{\deleted{\unsure{0}}}{\mathcal M}_{\rho,\sigma}^{\unsure{\flat|*|\natural}}
\end{tikzcd}}
\]
). Finally
\[
i_G(\deleted{\unsure{f}}\inserted{\unsure{x}})(i(\lambda))
=i_G(x)\cdot(i_B(\lambda)i_G(\lambda^{-1}))\cdot i_G(x)^{-1}
\]
now by definition, $i_B(\lambda)$ and $i_G(\lambda)$ operate
in the same way on $G$, hence
$i(\lambda)=i_B(\lambda)i_G(\lambda^{-1})$
operates trivially on it, hence the last factor
$i_G(x)^{-1}$ can \unsure{jump} \unsure{over|under} the factor
$i(\lambda)$, we find
\[
i_G(x)i_G(x)^{-1}i(\lambda)=i(\lambda)
\]
\margin{O.K.}
\margin{391}

% Printed archive footer (not autograph).
\begin{center}
391\\
\texttt{https://grothendieck.umontpellier.fr}\\
Université de Montpellier, UMR CNRS 5149, Case courrier 051, Place Eugène Bataillon\\
34095 Montpellier cedex 5 - France
\end{center}

\scanpage{393}

\uline{Remark} $\widetilde S\mathcal M_{\rho,\sigma}$ is a flat group of relative dimension
$\underbrace{2}_{B'_0}+\underbrace{2}_{T_\rho}+\underbrace{2}_{N'_\sigma}+\underbrace{4}_{G}=10$,
$S\mathcal M_{\rho,\sigma}$ is of relative dimension
$\underbrace{\unsure{4}}_{S\mathcal G_{\rho,\sigma}^{*}}+\underbrace{3}_{SG}=7$,
on passing to the quotient by $\underset{\simeq\mathbb G_a}{i(L_0)}$
\unsure{gives|to give} groups \deleted{\unsure{this}} denoted
$\widetilde{\mathcal M}_{\rho,\sigma}$, $\mathcal M_{\rho,\sigma}$,
of relative dimension $9$ resp. $6$, giving rise to a diagram
\[
(43)\qquad
\diagram{\begin{tikzcd}[ampersand replacement=\&,column sep=large,row sep=large]
1 \arrow[r] \& \mathbb G_a \arrow[r,"\lambda\mapsto\varepsilon_0(\lambda)"] \arrow[d,equal]
  \& S\mathcal M_{\rho,\sigma} \arrow[r] \arrow[d,hook,"\unsure{\mathrm{can.}}"']
  \& \mathcal M_{\rho,\sigma} \arrow[r] \arrow[d,hook] \& 1 \\
1 \arrow[r] \& \mathbb G_a \arrow[r]
  \& \widetilde S\mathcal M_{\rho,\sigma} \arrow[r]
  \& \unsure{\mathcal M_{\rho,\sigma}|\widetilde{\mathcal M}_{\rho,\sigma}} \arrow[r,no head] \& 1
\end{tikzcd}}
\]

The passage to the quotient presents no difficulties --- the quotients are in fact
\unsure{representable} \deleted{\unsure{by the|by}}\inserted{as} $1/2$ direct products
\[
(44)\qquad
\left\{\begin{aligned}
\deleted{\unsure{S}}\mathcal M_{\rho,\sigma}
 &\simeq \deleted{\unsure{T_0}}\,S\mathcal G_{\rho,\sigma}^{*}(0)\cdot SG\\
\widetilde{\mathcal M}_{\rho,\sigma}^{\deleted{\unsure{\sim}}}
 &\simeq \mathcal G_{\rho,\sigma}^{*\deleted{\unsure{\flat|\natural|\circ}}}(\unsure{0|\omega})\cdot\deleted{\unsure{S}}G
\end{aligned}\right.
\]
where the subgroup schemes
$S\mathcal G_{\rho,\sigma}^{*\unsure{\flat|\natural|\circ}}(\unsure{0|\omega})$ and
$\mathcal G_{\rho,\sigma}^{*\unsure{\flat|\natural|\circ}}(\unsure{0|\omega})$ of
$S\mathcal G_{\rho,\sigma}^{*}$, $\mathcal G_{\rho,\sigma}^{*}$
are defined in $\boxed{\text{\uline{VIII}}}$,
$\mathcal G_{\rho,\sigma}^{*\unsure{\flat|\natural|\circ}}(\unsure{0|\omega})$ being identified with
$T_0\times T_\rho\times N'_\sigma$, and
$S\mathcal G_{\rho,\sigma}^{*\unsure{\flat|\natural|\circ}}(0)$ with its intersection
with $S\mathcal G_{\rho,\sigma}^{*}$ (defined by the supplementary conditions $(20)$ on page $487'$).

\margin{392}

% Printed archive footer (not autograph).
\begin{center}
392\\
\texttt{https://grothendieck.umontpellier.fr}\\
Université de Montpellier, UMR CNRS 5149, Case courrier 051, Place Eugène Bataillon\\
34095 Montpellier cedex 5 - France
\end{center}

\scanpage{394}
\mspage{498}

Returning now to the arithmetic situation, we find
\inserted{by passage to the quotient via $(37)$} a canonical \unsure{homomorphism}
\margin{in view of $(13)$}
\[
(45)\qquad
\diagram{\begin{tikzcd}[ampersand replacement=\&,column sep=large]
\deleted{\unsure{S_0}}\mathcal M_{0,3}^{!\sim} \arrow[r]
\& \underbrace{\mathcal M_{\rho,\sigma}(\widehat{\mathbb Z})}_{\begin{gathered}
\simeq\\
\underbrace{\deleted{\unsure{B'_0}}\unsure{T_0|\widehat T_0}\times
T_\rho(\widehat{\mathbb Z})\times T_\sigma(\widehat{\mathbb Z})}_{\begin{gathered}
\text{operating on}\\
\mathrm{Sl}(2,\widehat{\mathbb Z})\ \text{via the}\\
\text{factor }T_0
\end{gathered}}\cdot\mathrm{Sl}(2,\widehat{\mathbb Z})
\end{gathered}}/\Sigma
\end{tikzcd}}
\]
$\Sigma$ being the subgroup $\simeq\{\pm1\}^3$, product
of three factors $\mathbb Z/2\mathbb Z$ included resp.$^{\unsure{t}}$
in $T_\rho(\widehat{\mathbb Z})$, $T_\sigma(\widehat{\mathbb Z})$
and in $\mathrm{Sl}(2,\widehat{\mathbb Z})$.

\uline{Remark} This homomorphism was obtained here,
\inserted{essentially,} with the help of the homomorphism from
$\mathcal M_{0,3}^{!\sim}\unsure{(0)}$
\[
(46)\qquad
\diagram{\begin{tikzcd}[ampersand replacement=\&,column sep=large,row sep=huge]
\begin{gathered}
\mathcal M_{0,3}^{!\sim}(0)\\
\text{subgroup}\\
\text{of the automorphisms}\\
\text{``normalized''}
\end{gathered} \arrow[dr]
\& \begin{gathered}
=\operatorname{Norm}_{\mathcal M_{0,3}^{!\sim}}(L_0)\\
\text{(cf. }\S46,\ \text{\uline{V}}\text{)}
\end{gathered} \arrow[dr,-{Circle[length=2pt,width=2pt]}] \& \\
\& S\mathcal G_{\rho,\sigma}^{*\unsure{\flat|\natural|\circ}}(0)(\widehat{\mathbb Z})/\Sigma_0
  \arrow[r,hook,no head]
\& S\mathcal G_{\rho,\sigma}^{*}(\widehat{\mathbb Z})/\Sigma_0
\end{tikzcd}}
\]
and of the composite
\[
\diagram{\begin{tikzcd}[ampersand replacement=\&,column sep=large,row sep=small]
\deleted{\unsure{\mathcal G}}\widehat\Pi_{0,3}
  \arrow[r,hook,no head]
\& \widehat{\unsure{\mathcal G|\mathcal C}}_{0,3}^{+}
  \arrow[r] \arrow[d,phantom,"\unsure{\cup|\wr}" description]
\& \mathrm{Sl}(2,\widehat{\mathbb Z})' \\
\& \mathrm{Sl}(2,\mathbb Z)^{\prime\wedge} \&
\end{tikzcd}}
\]
by writing
\[
(47)\qquad
\mathcal M_{0,3}^{!\sim}\simeq\mathcal M_{0,3}^{!\sim}\cdot\widehat\Pi_{0,3}\quad\unsure{!|:}
\]

\margin{393}

% Printed archive footer (not autograph).
\begin{center}
393\\
\texttt{https://grothendieck.umontpellier.fr}\\
Université de Montpellier, UMR CNRS 5149, Case courrier 051, Place Eugène Bataillon\\
34095 Montpellier cedex 5 - France
\end{center}

\scanpage{395}

But we also have a semidirect product
\[
(48)\qquad
\mathcal M_{0,3}^{\sim}\simeq
\mathcal M_{0,3}^{!\sim}(\unsure{0|\omega})\cdot
\unsure{\mathcal G|\mathcal C}_{0,3}^{+}
\]
which, combined directly with
\[
\diagram{\begin{tikzcd}[ampersand replacement=\&,column sep=large]
\widehat{\unsure{\mathcal G|\mathcal C}}_{0,3}^{+}\arrow[r]
\&\mathrm{Sl}(2,\widehat{\mathbb Z})'
\end{tikzcd}}
\]
gives a homomorphism
\[
(49)\qquad
\diagram{\begin{tikzcd}[ampersand replacement=\&,column sep=large]
\mathcal M_{0,3}^{\sim}\arrow[r]
\& \bigl(\mathcal M_{\rho,\sigma}(\widehat{\mathbb Z})
\simeq S\mathcal G_{\rho,\sigma}^{*}(\widehat{\mathbb Z})\cdot
\mathrm{Sl}(2,\widehat{\mathbb Z})\bigr)/\unsure{\Sigma}
\end{tikzcd}}
\]
which extends the homomorphism $(45)$, defined first on its normal
subgroup $\mathcal M_{0,3}^{!\sim}$ of index $6$.

\margin{394}

% Printed archive footer (not autograph).
\begin{center}
394\\
\texttt{https://grothendieck.umontpellier.fr}\\
Université de Montpellier, UMR CNRS 5149, Case courrier 051, Place Eugène Bataillon\\
34095 Montpellier cedex 5 - France
\end{center}

\scanpage{396}
\mspage{499}

$\overline{\text{\uline{X}}}$\quad Study of the group
$\mathcal M_{\rho,\sigma}$ and of $\mathcal M_{\rho,\sigma}\to\mathrm{Gl}(2)$

By construction of the semidirect product
\[
(1)\qquad
\diagram{\begin{tikzcd}[ampersand replacement=\&,column sep=small,row sep=huge]
\mathcal M_{\rho,\sigma}^{\sim} \arrow[r,phantom,"=" description]
\& \mathcal G_{\rho,\sigma}^{*\unsure{\flat|\natural|\circ}}(\unsure{\omega|0})\cdot G
  \arrow[r,phantom,"\simeq" description] \arrow[d,phantom,"\cup" description]
\& \underbrace{(T_0\times T_\rho\times N'_\sigma)}_{\begin{gathered}
\text{operating on }G\\ \text{via }T_0\hookrightarrow G
\end{gathered}}\cdot G \\
\mathcal M_{\rho,\sigma} \arrow[r,phantom,"=" description]
\& S\mathcal G_{\rho,\sigma}^{*\unsure{\flat|\natural|\circ}}(\unsure{\omega|0})\cdot SG \arrow[r,hook]
\& \underbrace{(T_\rho\times N'_\sigma)}_{\begin{gathered}
\text{operating on }G\\
\text{via }N'_\sigma\xrightarrow{\det}\mathbb G_m
\xleftarrow[\sim]{\det}T_0\hookrightarrow G
\end{gathered}}\cdot G
\end{tikzcd}}
\]
we have a group homomorphism
\[
(2)\qquad
\left\{\begin{gathered}
\diagram{\begin{tikzcd}[ampersand replacement=\&,column sep=large]
\mathcal M_{\rho,\sigma}^{\sim}\arrow[r,"\Theta"] \& G
\end{tikzcd}}\\
\Theta\bigl(i_B(x)i_{\deleted{\unsure{\rho|B}}\inserted{\rho}}(y)
\unsure{i_\sigma|i'_\sigma}(z)i_G(f)\bigr)=x.f
\end{gathered}\right.
\]
\[
(x\in\Gamma T_0,\quad y\in\Gamma T_\rho,\quad
z\in\Gamma\unsure{N'_\sigma},\quad f\in\Gamma G).
\]

The composite homomorphism
\[
(3)\qquad
\diagram{\begin{tikzcd}[ampersand replacement=\&,column sep=large,row sep=large]
S\mathcal M_{\rho,\sigma}^{\sim}\arrow[r,no head]
  \arrow[d,phantom,"\unsure{\wr}" description]
\&\mathcal M_{\rho,\sigma}^{\sim}\arrow[r,"\Theta"] \& G\\
\underbrace{\deleted{\unsure{\mathcal M}}\mathcal G_{\rho,\sigma}^{*\unsure{\flat|\natural|\circ}}}_{B'_0\times T_\rho\times N'_\sigma}\cdot G
\&\&
\end{tikzcd}}
\]
\margin{(composite still denoted $\Theta$)}
\deleted{is} described by
\[
\Theta\bigl(i_B(x)i_\rho(y)i_\sigma(z)i_G(\unsure{f|\varphi})\bigr)=x.f
\]
and $\Theta$ will indeed be trivial on the
\[
i(\lambda)=i_B(\lambda)i_G(\lambda^{-1})\qquad(\lambda\in\Gamma L_0)
\]
and therefore passes to the quotient as $\Theta:\mathcal M_{\rho,\sigma}^{\sim}\to G$.

\margin{395}

% Printed archive footer (not autograph).
\begin{center}
395\\
\texttt{https://grothendieck.umontpellier.fr}\\
Université de Montpellier, UMR CNRS 5149, Case courrier 051, Place Eugène Bataillon\\
34095 Montpellier cedex 5 - France
\end{center}
